\Qhead{6.4}{Линейный фильтр Калмана--Бьюси}{Линейный фильтр Калмана-Бьюси. Стохастические уравнения для модели сообщения и шума. Дифференциальные уравнения фильтра. Уравнение для апостериорной информации в форме уравнения Риккати. Сравнение фильтрации методом Колмогорова-Винера и Калмана-Бьюси.}
\Subfile{6.4.1}{sections/06_detection/q04/s01}

\recap{Фильтр соединяет динамический прогноз и инновацию. Усиление определяется относительной неопределённостью, ковариация подчиняется Риккати. В линейно-гауссовской модели это точное апостериорное среднее.}
\board{Два стохастических уравнения; инновация; вывод ковариации и усиления; скалярный стационарный корень.}
\oral{Почему ковариация фильтра зависит от модели и шума, но не обязана зависеть от конкретных значений гауссовских наблюдений?}
\bait{Можно оставить информационную форму: обращение $P$ превращает измерительный член в положительный источник точности $H^{\mathsf T}R^{-1}H$.}
\sources{Сёрккя и Солин~\cite{sarkka}, §10.4, печатные с.~208--210 (PDF-с.~216--218); Wunsch~\cite{wunsch}, §6.2, дискретная иллюстрация. Вывод ковариации и информационная форма выполнены непосредственно.}
